\section{Cumulants, effective channels, and the commutant test}
\label{supp:channel_theorem}

This appendix proves the finite-channel statements used in the main text and
fixes the complex covariance convention.  None of the arguments assumes a
large fiber rank, a thermodynamic limit, or an ETH ansatz.

\subsection{Realification and weighted cumulants}

Let $\mathcal V$ be a finite-dimensional complex vector space and
$\mathcal V_{\mathbb R}$ its realification.  Choose a real basis and denote
the real coordinate vector of $Y\in\mathcal V$ by $y_{\mathbb R}$.  For a
real dual vector $t$, define
\begin{equation}
 \Phi_Y(t)=\mathbb E\exp\!\left(i\langle t,y_{\mathbb R}\rangle\right),
 \qquad
 K_Y(t)=\log\Phi_Y(t).
 \label{supp:eq:cumulant_generator}
\end{equation}
In a neighborhood of the origin where the required derivatives exist, the
order-$k$ cumulant tensor is
\begin{equation}
 \kappa_k(Y)[t_1,\ldots,t_k]
 =\left.\frac{1}{i^k}
 \partial_{s_1}\cdots\partial_{s_k}
 K_Y\!\left(\sum_{r=1}^{k}s_rt_r\right)
 \right|_{s_1=\cdots=s_k=0} .
 \label{supp:eq:cumulant_definition}
\end{equation}
Finite $k$th absolute moments of the tested real-linear projections are
sufficient for the derivative statement used here.  This construction is
real multilinear.  Extending it complex linearly recovers cumulants written
with $Y$ and $Y^\dagger$ in any chosen order.

Suppose $Y_1,\ldots,Y_n$ are mutually independent and take values in the same
$\mathcal V$.  A fixed outcome-independent linear identification may be used
if the original channel spaces differ.  For deterministic real weights and
$Z=\sum_\alpha w_\alpha Y_\alpha$, independence gives
\begin{align}
 \Phi_Z(t)
 &=\prod_{\alpha=1}^{n}\Phi_{Y_\alpha}(w_\alpha t),\nonumber\\
 K_Z(t)
 &=\sum_{\alpha=1}^{n}K_{Y_\alpha}(w_\alpha t).
 \label{supp:eq:cgf_factorization}
\end{align}
Substitution of Eq.~\eqref{supp:eq:cgf_factorization} into
Eq.~\eqref{supp:eq:cumulant_definition} yields
\begin{equation}
 \kappa_k(Z)[t_1,\ldots,t_k]
 =\sum_{\alpha=1}^{n}w_\alpha^k
 \kappa_k(Y_\alpha)[t_1,\ldots,t_k].
 \label{supp:eq:weighted_cumulant_proof}
\end{equation}
Since Eq.~\eqref{supp:eq:weighted_cumulant_proof} holds for arbitrary test
fields, it proves the tensor identity quoted in the main text.  Centering is
required for the simple
second-moment formulas and normalized fourth statistics used in the paper;
the abstract additivity of cumulants itself also holds for noncentered
variables when their first cumulants are retained.

Mutual independence is essential.  Pairwise independence suffices only up to
orders controlled by the corresponding joint laws, and zero cross covariance
does not remove a mixed fourth cumulant.  If the channel weights depend on
the observed $Y_\alpha$, Eq.~\eqref{supp:eq:cgf_factorization} likewise does
not apply.

\subsection{Complete complex Wick subtraction}

For a centered complex vector $z$, our conventions are
\begin{equation}
 C_{ij}=\mathbb E[z_i z_j^*],
 \qquad
 \Pi_{ij}=\mathbb E[z_i z_j].
 \label{supp:eq:complex_second_moments}
\end{equation}
The covariance is Hermitian positive semidefinite and the pseudocovariance is
complex symmetric.  In terms of $z=u+iv$, the pair $(C,\Pi)$ contains the
same information as the full real covariance of $(u,v)$.  In particular,
$\Pi=0$ is a property of a proper complex population, not a convention that
can be imposed after sampling.

For the ordering $z_i z_j^*z_k z_l^*$, Isserlis contraction of a centered
complex Gaussian gives
\begin{equation}
 W_{i\bar j k\bar l}
 =C_{ij}C_{kl}+C_{il}C_{kj}+\Pi_{ik}\Pi_{jl}^*.
 \label{supp:eq:complex_wick_tensor}
\end{equation}
The corresponding connected tensor is
\begin{equation}
 \Gamma_{i\bar j k\bar l}
 =\mathbb E[z_i z_j^*z_k z_l^*]
 -W_{i\bar j k\bar l}.
 \label{supp:eq:complete_connected_tensor}
\end{equation}
Other index orderings follow from the same realified covariance and contain
the analogous three pairings.  The scalar contraction used by the executable
oracle is
\begin{align}
 \widehat\kappa_4^{\mathrm{conn}}
 &=\mathbb E|z|^4
 -2(\mathbb E|z|^2)^2
 -|\mathbb E z^2|^2,\nonumber\\
 \widehat R_4^{\mathrm{full}}
 &=\frac{\widehat\kappa_4^{\mathrm{conn}}}
 {(\mathbb E|z|^2)^2} .
 \label{supp:eq:scalar_complete_kappa}
\end{align}
Equation~\eqref{supp:eq:scalar_complete_kappa} displays explicitly why
ordinary covariance alone is not a complete Gaussian null.

For independent centered channel vectors $y_\alpha$, direct expansion gives
\begin{equation}
 C_Z=\sum_\alpha w_\alpha^2 C_\alpha,
 \qquad
 \Pi_Z=\sum_\alpha w_\alpha^2\Pi_\alpha.
 \label{supp:eq:weighted_cp}
\end{equation}
At fourth order, all terms involving two distinct independent channels are
exactly the Wick pairings constructed from Eq.~\eqref{supp:eq:weighted_cp}.
They cancel in Eq.~\eqref{supp:eq:complete_connected_tensor}, leaving only
$\sum_\alpha w_\alpha^4\Gamma_\alpha$, as required by the general cumulant
proof.

\subsection{Participation number and covariance eigenvalues}

Put
\begin{equation}
 q_\alpha=\frac{w_\alpha^2}{\sum_\beta w_\beta^2},
 \qquad \sum_\alpha q_\alpha=1 .
 \label{supp:eq:normalized_channel_power}
\end{equation}
Then
\begin{equation}
 N_{\mathrm{eff}}=\frac{1}{\sum_\alpha q_\alpha^2}.
 \label{supp:eq:participation_ratio}
\end{equation}
This is the inverse participation ratio of the channel powers, not a fitted
sample size.  If the independent, variance-normalized channels define
orthogonal directions in a channel-label covariance matrix
$G=\operatorname{diag}(q_1,\ldots,q_n)$, then
\begin{equation}
 N_{\mathrm{eff}}
 =\frac{(\operatorname{Tr}G)^2}{\operatorname{Tr}G^2}
 =\frac{(\sum_r g_r)^2}{\sum_r g_r^2},
 \label{supp:eq:covariance_eigen_neff}
\end{equation}
where $g_r$ are the eigenvalues of $G$.  This eigenvalue form explains the
usual covariance participation ratio.

Equation~\eqref{supp:eq:covariance_eigen_neff} must not be applied blindly to
the covariance of the flattened physical response.  Overlapping channel
tensors can have a response covariance whose eigenvalues do not identify the
independent units.  For correlated channel labels one can define the
descriptive quantity
$(\operatorname{Tr}R)^2/\operatorname{Tr}R^2$ from a label covariance $R$,
but the mixed connected tensors omitted from
Eq.~\eqref{supp:eq:weighted_cumulant_proof} remain.  Thus a covariance
participation ratio does not restore the theorem when independence is absent.

If every channel has the same variance-normalized connected tensor $K_4$,
Eq.~\eqref{supp:eq:weighted_cumulant_proof} and
Eq.~\eqref{supp:eq:participation_ratio} give
\begin{equation}
 \frac{\kappa_4(Z)}{(\sum_\alpha w_\alpha^2)^2}
 =K_4\sum_\alpha q_\alpha^2
 =\frac{K_4}{N_{\mathrm{eff}}}.
 \label{supp:eq:neff_derivation}
\end{equation}
If only $\|\kappa_4(Y_\alpha)\|\le K$ is known, the triangle inequality gives
the corresponding $K/N_{\mathrm{eff}}$ bound.  If the tensors are known but
unequal, they must remain inside the weighted sum; replacing them by one
average is an additional modeling assumption.

\subsection{Commutant as a linear system}

Let $A_1,\ldots,A_m$ be generators of $\mathcal A_X$, for example the
products $X_a^\dagger X_b$ together with the adjoints needed for star closure.
A matrix $M$ lies in the commutant precisely when
\begin{equation}
 MA_r-A_rM=0,
 \qquad r=1,\ldots,m .
 \label{supp:eq:commutator_equations}
\end{equation}
With column-major vectorization,
\begin{equation}
 \operatorname{vec}(MA_r-A_rM)
 =\left(A_r^T\otimes I-I\otimes A_r\right)
 \operatorname{vec}M .
 \label{supp:eq:commutator_superoperator}
\end{equation}
Stacking the superoperators in
Eq.~\eqref{supp:eq:commutator_superoperator} gives a finite linear system.
Its complex nullity is $\dim\mathcal A_X'$.  The executable audit evaluates
that nullity by a singular-value decomposition at a stated tolerance.  For
two Pauli generators it returns one; for one nondegenerate diagonal generator
it returns two; for the identity alone it returns four.

For completeness, suppose $\mathcal A_X$ is a finite-dimensional complex
star algebra and has a nontrivial invariant subspace.  Star closure makes its
orthogonal complement invariant as well, and the orthogonal projection onto
the subspace is a nonscalar member of $\mathcal A_X'$.  Hence a scalar
commutant implies irreducibility.  Burnside's theorem then identifies an
irreducible unital complex matrix algebra with the full endomorphism algebra.
This argument is algebraic.  It contains no distribution, cumulant, or
large-size statement and therefore supplies no implication from
irreducibility to ETH.

\subsection{Executable oracle and claim boundary}

The synthetic channel population uses independent entries
\begin{equation}
 y=\frac{B\,\xi}{\sqrt p},
 \qquad
 \Pr(B=1)=p=0.1,
 \qquad
 \xi\in\{1,i,-1,-i\}
 \label{supp:eq:oracle_population}
\end{equation}
with the four phases equally likely.  It has
$\mathbb Ey=0$, $\mathbb E|y|^2=1$, $\mathbb Ey^2=0$, and
\begin{equation}
 \kappa_4(y)=\mathbb E|y|^4-2=\frac{1}{p}-2=8 .
 \label{supp:eq:oracle_kappa}
\end{equation}
The audit draws $50{,}000$ complex $2\times2$ matrices, uses $25$ blocks for
uncertainty, and evaluates equal and linearly unequal weights for
$n=4,8,16,32,64$.  Its maximum relative discrepancy from the fixed
Eq.~\eqref{supp:eq:neff_derivation} prediction is $\ChannelTheoryMaximumRelativeError$; the
registered implementation threshold is $0.08$.  The exponent is not fitted.

All pass/fail decisions use the unrounded double-precision calculation.  A
recursive ten-significant-digit quantization is applied only when serializing
the JSON record and computing its portable scientific hash; nonfinite values
are rejected.  This serialization convention has no effect on the numerical
gate.

The derivation fails or changes in any of the following cases: mutually
correlated channels, absent fourth moments, outcome-dependent weights,
unresolved symmetry blocks, a size-dependent tangent population, a dominant
weight that prevents $N_{\mathrm{eff}}$ from growing, or a closing external
gap.  These are hypotheses to test, not technical qualifications that may be
discarded after observing a mismatch.  In the prospective chiral calculation
the fixed random-channel validation fails, so the synthetic oracle and exact
proof cannot be promoted to a physical Geometric-ETH conclusion.
